\chapter*{Abstract}
\phantomsection
\addcontentsline{toc}{chapter}{Abstract}

Scalar-on-function linear regression is an inverse problem in which the
predictor covariance operator determines both identification and the
stability of slope recovery.  This thesis examines how regularisation,
numerical representation, and the purpose of an analysis interact.  It
compares functional principal component regression, raw and Arnoldi
implementations of functional partial least squares, and a
conjugate-gradient-type procedure based on moment-residual minimisation.

The theoretical discussion distinguishes ordinary slope error,
covariance-weighted prediction error, and the residual approximation needed
for inference.  Monte Carlo experiments with 5,000 replications in each of
three controlled designs show that the implemented early-stopping rule can
reduce slope variance while leaving predictive signal unresolved.  Raw and
orthogonalised Krylov fits agree closely when their coordinates are reliable,
but rare ill-conditioned raw fits create severe error tails.  A separate
inference study shows that an estimation-oriented stopping variant can leave
a substantial moment residual on the root-$n$ scale.  Late iteration makes this
algorithmic discrepancy negligible in the studied designs; calibration
remains conditional on the reference distribution, discretisation, and
finite-sample approximation.

An application uses instrument-1 near-infrared spectra to predict assay for
pharmaceutical tablets, with 195 development observations and a designated
460-tablet benchmark.  Raw and Arnoldi FPLS achieve benchmark RMSEs of 4.607
and 4.643 in the supplied assay units.  Their predictions are very close,
while Arnoldi uses a stable five-dimensional representation and the selected
raw representation is severely ill-conditioned.  The benchmark had been
inspected during dataset selection, and the uncertainty intervals condition
on the fitted models.  The findings support matching regularisation and
stopping to the statistical target, monitoring numerical reliability, and
distinguishing controlled inferential benchmarks from feasible data analysis.
